% =====================================================================
%  3.4 Rasjonale funksjoner — fasit
% =====================================================================

\fk{Del 1 — Asymptoter og hull}
\fo{1.1} a) $x = 2$, $y = 1$ \ b) $x = -3$, $y = 2$ \ c) $x = -1$, $y = 0$
\fo{1.2} a) $4$ \ b) Brøken nærmer seg 0. \ c) $y = 3$
\fo{1.3} a) $D_f = \{x \mid x \neq -4\}$ \ b) $x = 3$ \ c) $\left(0, -\frac{3}{2}\right)$
\fo{1.4} a) $(5x - 2) : (x - 3) = 5$ rest $13$, så $f(x) = 5 + \frac{13}{x-3}$.
b) Vertikal: $x = 3$. Horisontal: $y = 5$.
c) Nullpunkt $x = \frac{2}{5}$. Skjæring med $y$-aksen: $\left(0, \frac{2}{3}\right)$.
d) To grener: den høyre ligger over $y = 5$ og går mot asymptotene. Den venstre ligger under $y = 5$. Grafen går gjennom $\left(\frac{2}{5}, 0\right)$ og $\left(0, \frac{2}{3}\right)$.
\fo{1.5} a) $p(x) = \frac{(x-1)(x-4)}{x-1} = x - 4$ for $x \neq 1$. Hull i $(1, -3)$. Ingen asymptoter.
b) $q(x) = \frac{(x-1)(x-2)}{(x-2)(x+2)} = \frac{x-1}{x+2}$ for $x \neq \pm 2$. Hull i $\left(2, \frac{1}{4}\right)$. Vertikal asymptote $x = -2$. Horisontal asymptote $y = 1$.
\fo{1.6} Den første feilen er tredje trinn: $x = 4$ er ikke en asymptote. Faktoren $x - 4$ ble forkortet bort, så grafen har et hull i $\left(4, \frac{1}{4}\right)$.
Riktig: vertikal asymptote $x = -4$, hull i $\left(4, \frac{1}{4}\right)$, horisontal asymptote $y = 0$.
\fo{1.7} Ett svar: $f(x) = \frac{-2x+2}{x-3}$. Nevneren er 0 for $x = 3$, og telleren er $-4 \neq 0$. Horisontal asymptote $\frac{-2}{1} = -2$. Telleren er 0 for $x = 1$.
Flere svar er riktige, for eksempel $\frac{-4x+4}{2x-6}$.
\fo{1.8} a) Noen ganger. Er $ad \neq bc$, er det en vertikal asymptote $x = -\frac{d}{c}$. Er $ad = bc$, forkortes brøken til en konstant. For eksempel er $\frac{2x+4}{x+2} = 2$ for $x \neq -2$, med et hull og uten asymptote.
b) Aldri. $f(x) - \frac{a}{c} = \frac{bc - ad}{c(cx+d)}$. Telleren er et tall som ikke er 0, så brøken er aldri 0.
c) $\frac{3x^2+x}{x^2+1} = 3$ gir $3x^2 + x = 3x^2 + 3$, altså $x = 3$. Punktet $(3, 3)$ ligger på grafen. Påstanden stemmer ikke generelt, men for $\frac{ax+b}{cx+d}$ stemmer den (se b).
\fo{1.9} Ett svar: $f(x) = \frac{15x(x+1)}{(x+1)(x-4)}$. Etter forkorting er $f(x) = \frac{15x}{x-4}$. Vertikal asymptote $x = 4$. Hull ved $x = -1$ med $f(-1) = \frac{-15}{-5} = 3$. Flere svar er riktige.
\fo{Sjekk 1} a) $x = 4$ og $y = 6$ \ b) Mot en asymptote nærmer grafen seg uten å nå linjen. Et hull er ett enkelt punkt som mangler, der teller og nevner begge er 0.
c) $(x+4) : (x-2) = 1$ rest $6$. Når $x$ blir stor, går $\frac{6}{x-2}$ mot 0, så $f(x)$ nærmer seg 1.

\fk{Del 2 — Rasjonale likninger og ulikheter}
\fo{2.1} a) $x = 3$ \ b) $x = 2$ \ c) $x = 3$
\fo{2.2} a) $x \neq 3$ \ b) $x \neq 0$ og $x \neq -4$ \ c) $x \neq 3$ og $x \neq -3$
\fo{2.3} a) Produktlinja er heltrukket for $x < -3$ og $x > 1$, og stiplet mellom. Den har $\times$ i $-3$ og $0$ i $1$.
b) $\langle \leftarrow, -3\rangle \cup \langle 1, \rightarrow\rangle$
\fo{2.4} a) Vilkår $x \neq 3$. Fra $2x = x - 3 + 6$ får vi $x = 3$. Det er forbudt, så ingen løsning.
b) Vilkår $x \neq 0, -2$. $4(x+2) + 4x = 3x(x+2)$ gir $3x^2 - 2x - 8 = 0$. Løsninger $x = 2$ og $x = -\frac{4}{3}$.
c) Vilkår $x \neq 1, -2$. $2(x+2) = x(x-1)$ gir $x^2 - 3x - 4 = 0$. Løsninger $x = 4$ og $x = -1$.
\fo{2.5} a) $\langle \leftarrow, -1\rangle \cup [4, \rightarrow\rangle$
b) $\frac{2x-1}{x-3} - 1 = \frac{x+2}{x-3} < 0$ gir $\langle -2, 3\rangle$
c) $\frac{3}{x+2} - 1 = \frac{1-x}{x+2} \le 0$ gir $\langle \leftarrow, -2\rangle \cup [1, \rightarrow\rangle$
\fo{2.6} Den første feilen er overgangen til $\frac{x+5-2x-2}{x-1} \ge 0$. Det skal være $-2(x-1) = -2x + 2$, ikke $-2x - 2$.
Riktig: $\frac{x+5-2x+2}{x-1} = \frac{7-x}{x-1} \ge 0$. Svar: $\langle 1, 7]$.
\fo{2.7} Brøken er 0 i $x = 3$ og ikke definert i $x = -2$. Fortegnet skifter i begge. Ulikheten er $\frac{x-3}{x+2} \ge 0$, så $p = -3$ og $q = 2$. Nevnerens nullpunkt $-2$ er ikke med, og telleren 3 er med.
\fo{2.8} a) $\langle -3, 1\rangle \cup \langle 1, \rightarrow\rangle$ \ b) $\langle -3, \rightarrow\rangle$
Telleren $(x-1)^2$ er aldri negativ og er 0 bare for $x = 1$. Fortegnet skifter ikke der. Brøken er derfor positiv overalt der $x + 3 > 0$, bortsett fra i $x = 1$, der den er 0. Med $>$ er 0 ikke en løsning. Med $\ge$ er 0 en løsning.
\fo{2.9} $x + k = 3(x - 2)$ gir $x = \frac{k+6}{2}$. Løsningen er forbudt når $x = 2$, altså $k = -2$. Da er $\frac{x-2}{x-2} = 1 \neq 3$ for alle tillatte $x$. Svar: $k = -2$.
\fo{Sjekk 2} a) $x = 3$ \ b) $\langle -3, 2\rangle$ \ c) Vi vet ikke om $x$ er positiv eller negativ. For $x < 0$ må tegnet snus. Riktig løsning er $0 < x < \frac{5}{2}$.

\fk{Ekstra utfordring}
\fo{U1} a) $(x^2 + 1) : (x - 1) = x + 1$ rest $2$, så $f(x) = x + 1 + \frac{2}{x-1}$.
b) Brøkdelen $\frac{2}{x-1}$ går mot 0. Grafen nærmer seg den skrå linja $y = x + 1$.
c) Vertikal asymptote $x = 1$. Ingen horisontal asymptote, for $f(x)$ vokser som $x + 1$.
\fo{U2} a) $K(x) = 20 + \frac{5000}{x}$. Horisontal asymptote $y = 20$: kostnaden per enhet nærmer seg 20 kr, som er kostnaden per enhet uten faste kostnader.
b) $\frac{5000 + 20x}{x} < 25$ gir $5000 + 20x < 25x$, altså $x > 1000$. (Her er $x > 0$, så vi kan gange med $x$.)
c) Nei, $K(0)$ finnes ikke. Modellen gir at kostnaden aldri blir under 20 kr, også for veldig store $x$. Det er urealistisk hvis bedriften må utvide produksjonen.
\fo{U3} a) \lstinline[style=laereverk,numbers=none,frame=none]|return (2 * x + 1) / (x - 1)|. Utskrift: 1.1 32.0, 1.01 302.0, 1.001 3002.0, 10 2.333, 100 2.03, 1000 2.003.
b) Når $x$ nærmer seg 1, blir $f(x)$ veldig stor (vertikal asymptote $x = 1$). Når $x$ blir stor, nærmer $f(x)$ seg 2 (horisontal asymptote $y = 2$).
c) $f(1)$ krever deling på 0. Python gir feilmeldingen \texttt{ZeroDivisionError}.
\fo{U4} $\frac{1}{x-1} - \frac{1}{x+2} = \frac{3}{(x-1)(x+2)} \ge 0$. Telleren 3 er positiv, så produktet $(x-1)(x+2)$ må være positivt. Svar: $\langle \leftarrow, -2\rangle \cup \langle 1, \rightarrow\rangle$.
\fo{U5} a) $f(x) = 3 + \frac{7}{x-2}$. Brøken $\frac{7}{x-2}$ er aldri 0, så $f(x) \neq 3$.
b) $y(x-2) = 3x + 1$ gir $x(y-3) = 2y + 1$, altså $x = \frac{2y+1}{y-3}$ for $y \neq 3$. For $y = 3$ finnes ingen $x$.
c) $V_f = \{y \mid y \neq 3\}$. Generelt er $V_f = \{y \mid y \neq \frac{a}{c}\}$, for $\frac{ax+b}{cx+d} - \frac{a}{c} = \frac{bc-ad}{c(cx+d)}$ er aldri 0.
\fo{U6} $\frac{x+7}{x-1} = 1 + \frac{8}{x-1}$. Brøken er et helt tall nøyaktig når $x - 1$ går opp i 8, altså $x - 1 \in \{\pm1, \pm2, \pm4, \pm8\}$. Svar: $x \in \{-7, -3, -1, 0, 2, 3, 5, 9\}$.
\fo{U7} Ett svar: $f(x) = \frac{x-3}{(x+1)(x-2)(x-3)}$. Etter forkorting er $f(x) = \frac{1}{(x+1)(x-2)}$ for $x \neq 3$. Vertikale asymptoter $x = -1$ og $x = 2$. Grad i telleren er mindre enn i nevneren, så $y = 0$. Hullet er i $\left(3, \frac{1}{4}\right)$, fordi $\frac{1}{4 \cdot 1} = \frac{1}{4}$.
